\begin{abstract}
Weight-space model merging has emerged as a computationally efficient paradigm for consolidating multi-task capabilities, but dynamically routing inputs to task-specific experts on-the-fly remains a fundamental challenge under extreme data scarcity (e.g., $B_{\text{cal}} \le 64$). In such regimes, parametric routers suffer from severe overfitting, causing catastrophic generalization collapse on out-of-distribution (OOD) tasks. This work rejects current ad-hoc, heuristic regularizations in favor of a mathematically rigorous, first-principles approach. By formalizing dynamic weight-space merging, we prove that the Rademacher complexity of the dynamically merged hypothesis class is upper-bounded by the norms of the routing parameters weighted by the parameter-space magnitudes of the expert task vectors. Guided by this generalization bound, we introduce Spectral and Rademacher-guided Routing Regularization (SR3), an asymmetric regularization framework that penalizes the routing weights of each expert proportionally to the linear norm (Frobenius or Spectral) of its task vector. We also propose a smoothed $L_1$ Group-Lasso variant (SR3-L1) that directly minimizes the derived complexity bound. On a continuous weight-merging simulator featuring structured geometries and representation entanglement, we show that non-parametric methods collapse catastrophically, while parametric routing with SR3 achieves highly competitive joint multi-task accuracy comparable to state-of-the-art heuristics, demonstrating that aligning regularization intensity with parameter-space geometry is a theoretically sound and robust strategy for low-data generalization.
\end{abstract}